% Generated draft: requires mathematical and source review.
\documentclass[11pt]{article}
\usepackage[margin=25mm]{geometry}
\usepackage{fontspec}
\setmainfont{Latin Modern Roman}
\usepackage{amsmath,amssymb,mathtools,unicode-math}
\setmathfont{Latin Modern Math}
\usepackage{enumitem,needspace,xurl,hyperref}
\setlength{\parindent}{0pt}
\setlength{\parskip}{5pt}
\setlength{\emergencystretch}{3em}
\newcommand{\N}{\mathbb{N}}
\newcommand{\Z}{\mathbb{Z}}
\newcommand{\Q}{\mathbb{Q}}
\newcommand{\R}{\mathbb{R}}
\newcommand{\C}{\mathbb{C}}
\newcommand{\F}{\mathbb{F}}
\newcommand{\A}{\mathbb{A}}
\newcommand{\PP}{\mathbb{P}}
\newcommand{\E}{\mathbb{E}}
\newcommand{\eps}{\varepsilon}
\newcommand{\dd}{\,\mathrm{d}}
\begin{document}
\section*{UnsolvedMath ID: 20000230}
\textbf{Source problem code:} catalogue entry
\subsection{Definitions and mathematical statement}
The exact extracted problem begins ``Prove this for all $m$,'' but omits the antecedent map and hypotheses. Its binomial coefficient is damaged, and its discussion of ``$m$-wise sums'' is not defined. The verification says earlier PDF context recovers the setup, but that recovered context is not actually included in the supplied record. A precise unconditional all-$m$ theorem therefore cannot be reconstructed here.

A conditional interpretation consistent with the research report is the maximal-rank problem for section multiplication. For a smooth projective curve $X$ and line bundle $L$, put
\[
 \mu_m:\operatorname{Sym}^m H^0(X,L)\longrightarrow H^0(X,L^{\otimes m}).
\]
If $h^0(X,L)=r+1$, maximal rank means
\[
 \operatorname{rank}\mu_m
   =\min\left\{\binom{r+m}{m},h^0(X,L^{\otimes m})\right\}.
\]
The intended field, generality, genus, degree, embedding, and range of $m$ must be supplied before this is identified with ``this.'' The source says the existing $m=2$ argument is more natural in the injective case because no choice of independent sections is needed, describes a combinatorial game, and suggests higher syzygies afterward; these motivations do not supply the missing hypotheses.

The research aid invokes Larson's maximal rank theorem for general complex curves and reports a transition strip with
\[
 g=\binom r2+a,\qquad d=g+r,
\]
where the quadratic cokernel has dimension $a$ and the cubic kernel has dimension $\binom{r+1}{3}-2a$. The strip's actual range and line-bundle/general-embedding assumptions are not given; in particular the dimensions cannot be asserted for arbitrary $a$.

At $a=0$, it reports an arithmetically Cohen--Macaulay homogeneous coordinate ring, an ideal with pure $3$-linear resolution, and Betti formula
\[
 \beta_i=\binom{r+1}{i+2}\binom{i+1}{2},\qquad
 g=\binom r2,\quad d=\binom{r+1}{2}.
\]
The index convention and full resolution are not supplied. It proposes the union of all coordinate lines in $\PP^r$ as a characteristic-free monomial certificate with the same Hilbert series and Betti table, whose cubic basis is exactly the monomials supported on at most two variables. For $r=3$, it claims a flat Hilbert--Burch determinantal smoothing preserving the sixteen selected cubic monomials. These are unreviewed special-wall and certificate statements, not a replacement of the original all-$m$ question by quadratic/cubic cases alone.
\subsection{Short English statement}
The original all-$m$ assertion lacks its preceding definition and has damaged combinatorial notation. A multiplication-map maximal-rank interpretation is plausible from the report but remains conditional. The report supplies a complex-curve transition strip, a pure-resolution wall, and a rank-three monomial smoothing example, without the complete original hypotheses.
\subsection{Sources}
\begin{itemize}
\item Local UnsolvedMath registry, permanent ID $20000230$, source code \texttt{catalogue entry}; created and updated 2026-08-14. Status \texttt{solved}; classification \texttt{SOLVED-IN-LITERATURE}; difficulty L3. Proposer field null, while verification says preceding text suggests the problem by David Jensen.
\item AIM Problem Lists, Algebraic geometry, workshop \emph{Degenerations in algebraic geometry}, Problem 8.1, source type PDF, \url{https://aimath.org/pastworkshops/degenalggeomproblems.pdf}. Canonical location \texttt{aim-algebraic-geometry-notes.json}, \texttt{notes[229]}, tag \texttt{problem}; statement status \texttt{exact}. Verification explicitly notes omitted preceding map, damaged binomial coefficient, run-together words, and a reported page-4 inspection, but does not include its promised recovered context. Original remarks and literature field empty.
\item The research summary invokes Larson's maximal rank theorem and a Hilbert--Burch deformation without exact primary citations or the smoothing matrix.
\item Supplied research attempt 1: result \texttt{partial\_result}, run status \texttt{solved}, rigor label \texttt{proved}, low novelty confidence. Candidate contribution: derived wall theorem, Betti formula, characteristic-free coordinate-line certificate, and $r=3$ flat smoothing. Explicit caveat: machine-generated research aid, not peer reviewed; independently verify before citation.
\item Dataset attribution: Ulam AI / UnsolvedMath Contributors, CC BY 4.0; original problem attribution remains to the supplied AIM workshop collection.
\end{itemize}
\end{document}
